\documentclass[11pt,a4paper]{article}

\usepackage[T1]{fontenc}
\usepackage{lmodern}
\usepackage[margin=2.5cm]{geometry}
\usepackage{amsmath}
\usepackage{booktabs}
\usepackage{graphicx}
\usepackage{subcaption}
\usepackage[round]{natbib}
\usepackage[hidelinks]{hyperref}

\graphicspath{{figures/}}
\setlength{\emergencystretch}{1.5em}

\usepackage{url}
\DeclareUrlCommand\file{\urlstyle{tt}}  % file names, breakable at / and _
\newcommand{\eref}{^{\mathit{ref}}}

\title{Hardening Soil UMAT --- Test Report}
\author{ConstitutiveModels}
\date{October 5, 2026}

\begin{document}
\maketitle

% ---------------------------------------------------------------------------
\section{Scope}

This report summarises the verification of the Hardening Soil (HS) UMAT
\file{c_models/hardening_soil/hardening_soil.c}, following \citet{Schanz_1999}, \emph{The
hardening soil model: Formulation and verification}. Equation, section, table and figure numbers
without further reference refer to \citet{Schanz_1999}. The unit tests are in \file{tests/test_hardening_soil.py}; the figures are produced by
\file{docs/hs_test_report/make_figures.py}, which drives the UMAT along element test paths and
controls lateral stresses by Newton iteration on the strain increment where needed. The library
under test is \file{build_C/lib/hardening_soil.dll}, built from the current sources with
\file{Makefile.windows}.

The model combines
\begin{itemize}
  \item stress-dependent stiffness $E_{50}, E_{ur} \propto \left((\sigma_3 + a)/(p\eref + a)\right)^m$
        with $a = c\cot\varphi$ (Eqs.~3--5);
  \item shear hardening on the hyperbolic yield surface $f_{13}$ (Eqs.~7--9), after
        \citet{Kondner_1963} and \citet{Duncan_1970}, with Mohr--Coulomb failure at $q_f$ (Eq.~2);
  \item non-associated flow with the mobilised dilatancy angle of Rowe's stress--dilatancy theory
        \citep{Rowe_1962, Schanz_1996} (Eqs.~11--15), bounded so that the shear
        surface does not contract, and a dilatancy cut-off (Eqs.~38--39);
  \item an elliptic cap with associated flow and hardening of the pre-consolidation stress $p_c$
        (Eqs.~27--36);
  \item a return mapping in principal stress space with Koiter's rule at the triaxial corners
        \citep{Koiter_1960} (Eq.~26), and automatic sub-stepping.
\end{itemize}

Stresses and strains are compression positive, as in \citet{Schanz_1999}; the UMAT interface is tension
positive (Abaqus convention) and the drivers convert. $p = \frac{1}{3}\operatorname{tr}\boldsymbol{\sigma}'$,
$q = \sqrt{\frac{3}{2}\,\mathbf{s}:\mathbf{s}}$ (equal to $\sigma_1 - \sigma_3$ under triaxial
conditions), and $\gamma^p = \varepsilon_1^p - \varepsilon_2^p - \varepsilon_3^p$ (Eq.~9). Two
parameter sets are used (Table~\ref{tab:parameters}): the set of the unit tests, and loose Hostun
sand from Table~1 of \citet{Schanz_1999}.

\begin{table}[h]
  \centering
  \caption{Material parameters. $R_f = 0.9$ for Hostun sand is the default of Sec.~2 of \citet{Schanz_1999};
           its $M$ and $K_s/K_c$ are calibrated in Sec.~\ref{sec:hostun-oedometer}. Variations (cap
           off, $m = 0$, dilatancy cut-off, $\psi = 14^\circ$ for Hostun sand) are stated where they
           are used.}
  \label{tab:parameters}
  \small
  \begin{tabular}{lccccccccccc}
    \toprule
    Set & $E_{50}\eref$ & $E_{ur}\eref$ & $m$ & $c$ & $\varphi$ & $\psi$ & $p\eref$ & $R_f$ & $\nu_{ur}$ & $M$ & $K_s/K_c$ \\
        & [MPa] & [MPa] & [--] & [kPa] & [$^\circ$] & [$^\circ$] & [kPa] & [--] & [--] & [--] & [--] \\
    \midrule
    Unit tests   & 30 & 90 & 0.55 & 0 & 42 & 16 & 100 & 0.9 & 0.25 & 1.0   & 1.84  \\
    Hostun sand  & 20 & 60 & 0.65 & 0 & 34 & 0  & 100 & 0.9 & 0.20 & 1.233 & 2.109 \\
    \bottomrule
  \end{tabular}
\end{table}

% ---------------------------------------------------------------------------
\section{Test results}

All 18 unit tests pass (pytest, October 5, 2026).

\begin{table}[h]
  \centering
  \small
  \begin{tabular}{llc}
    \toprule
    Test & Checks & Result \\
    \midrule
    \texttt{elastic\_isotropic\_response}                & $E_{ur}$-based response inside the cap      & pass \\
    \texttt{shear\_yields\_immediately}                  & plastic shear from $q = 0$, $\gamma^p > 0$ & pass \\
    \texttt{hardening\_monotonic\_and\_yield\_...}       & $q$, $\gamma^p$ increase; $f_{13} \le 0$; $q \le q_f$ & pass \\
    \texttt{triaxial\_reaches\_and\_respects\_failure}   & drained triaxial, $q \le q_f$              & pass \\
    \texttt{principal\_frame\_invariance} ($\times 2$)   & rotated loading, rotated stress            & pass \\
    \texttt{drained\_triaxial\_follows\_hyperbola}       & Eq.~1; $E_{50}$ secant at $q_f/2$          & pass \\
    \texttt{plastic\_strains\_follow\_rowe\_dilatancy}   & Eqs.~9 and 11--15, no contraction          & pass \\
    \texttt{undrained\_shear\_surface\_...} ($\times 2$) & undrained, cap off: $p' = \sigma_c$ below $\varphi_{cv}$ & pass \\
    \texttt{cap\_isotropic\_normal\_compression}         & Eqs.~31--33; elastic unloading with $K_s$  & pass \\
    \texttt{mohr\_coulomb\_failure\_general\_...} ($\times 2$) & plane strain, triaxial extension     & pass \\
    \texttt{umat\_interface\_is\_tension\_positive}      & sign convention, regular tangent           & pass \\
    \texttt{tangent\_matches\_numerical\_...} ($\times 2$) & DDSDDE vs.\ finite differences           & pass \\
    \texttt{zero\_confinement\_and\_tension}             & apex return, positive definite tangent     & pass \\
    \texttt{stress\_update\_is\_continuous\_...}         & no jumps from the sub-stepping             & pass \\
    \bottomrule
  \end{tabular}
\end{table}

The tangent tests compare DDSDDE, the elasto-plastic (Koiter) tangent of the active yield
surfaces, with a finite-difference derivative of the stress update under continued plastic
loading, in a rotated frame:
\begin{itemize}
  \item \emph{cap off, plane strain} (no triaxial corner): relative difference
        $\|\mathbf{D} - \mathbf{D}_{FD}\|/\|\mathbf{D}_{FD}\| = 6.5\times10^{-4}$ (tolerance
        $10^{-2}$);
  \item \emph{cap on, triaxial compression corner}: $5.4\times10^{-3}$ (tolerance
        $3\times10^{-2}$). At the corner the exact tangent has no stiffness in the corner mode; the
        UMAT keeps 1\,\% of it (\file{HS_CORNER_STIFFNESS_FRACTION}) so that the global stiffness
        matrix stays regular, which accounts for most of the difference.
\end{itemize}
The remaining difference is the explicit (Euler) integration within the sub-steps: the stiffness
and dilatancy are frozen during a sub-step, while the tangent is evaluated at its end.

% ---------------------------------------------------------------------------
\section{Shear hardening: drained triaxial compression}

Drained triaxial compression at constant $\sigma_3$ uses the unit-test parameters, with the cap off
and no plastic volume change on the shear surface (dilatancy cut-off active from the start,
$e_0 = 0.8 > e_{cv} = 0.5$). Then $\gamma^p = 2\varepsilon_1^p$ and the hyperbola holds exactly. The
UMAT uses $E_i = 2E_{50}/(2 - R_f)$, which makes $E_{50}$ the secant stiffness at $q_f/2$ as defined
in Sec.~2.1; for $R_f \to 1$ this becomes the $E_i = 2E_{50}$ of Eq.~1.

The computed curves coincide with the hyperbola and stay below $q_f$ (Fig.~\ref{fig:hyperbola}a). The
largest relative deviation in $\varepsilon_1$ is $3.1\times10^{-4}$, at $\sigma_3 = 100$~kPa. The
secant stiffness at $q_f/2$ follows Eq.~3 within $1.1\times10^{-5}$ for $\sigma_3 = 25$--$800$~kPa
(Fig.~\ref{fig:hyperbola}b).

\begin{figure}[h]
  \centering
  \begin{subfigure}[t]{0.48\textwidth}
    \centering
    \includegraphics[width=\textwidth]{triaxial_hyperbola.pdf}
    \caption{$\varepsilon_1$--$q$ at three confining pressures}
  \end{subfigure}\hfill
  \begin{subfigure}[t]{0.48\textwidth}
    \centering
    \includegraphics[width=\textwidth]{triaxial_e50.pdf}
    \caption{secant stiffness $E_{50}$ vs.\ $\sigma_3$}
  \end{subfigure}
  \caption{Drained triaxial compression without cap and without plastic volume change (unit-test
           parameters).}
  \label{fig:hyperbola}
\end{figure}

% ---------------------------------------------------------------------------
\section{Stress dilatancy and dilatancy cut-off}

This case is drained triaxial compression at $\sigma_3 = 100$~kPa with $m = 0$ and the cap off. The
elastic stiffness is then constant, so the plastic strains follow directly from the total strains.
The plastic shear strain $\varepsilon_1^p - \varepsilon_2^p - \varepsilon_3^p$ equals the hardening
parameter $\gamma^p$ (Eq.~9) to $5\times10^{-14}$. The ratio
$-\Delta\varepsilon_v^p/\Delta\gamma^p$ gives the mobilised dilatancy angle.

Rowe's relation, Eq.~11, makes the soil contract below $\varphi_{cv}$, down to $\psi_m =
-\varphi_{cv}$ at $\varphi_m = 0$ (dotted line in Fig.~\ref{fig:dilatancy}a). The UMAT bounds it so
that the shear surface does not contract:
\begin{equation*}
  \sin\psi_m =
  \begin{cases}
    0 & \sin\varphi_m < \tfrac{3}{4}\sin\varphi \\[0.5ex]
    \max\left(\dfrac{\sin\varphi_m - \sin\varphi_{cv}}{1 - \sin\varphi_m\sin\varphi_{cv}},\, 0\right)
      & \sin\varphi_m \ge \tfrac{3}{4}\sin\varphi,\ \psi > 0 \\[1.5ex]
    \sin\psi & \sin\varphi_m \ge \tfrac{3}{4}\sin\varphi,\ \psi \le 0
  \end{cases}
\end{equation*}
Plastic compaction then comes from the cap alone. The computed curves of \citet{Schanz_1999} for loose Hostun
sand show no contraction of the shear surface either
(Secs.~\ref{sec:hostun-drained}--\ref{sec:undrained}).

Here $\varphi_{cv} = 28.85^\circ$ (Eq.~13), below the bound at $\varphi_m = 30.12^\circ$. The
computed $\sin\psi_m$ is zero (to $7\times10^{-15}$) below the bound. Above it the soil dilates
following Eq.~11, and reaches $\psi_m = \psi = 16.000^\circ$ at failure. The UMAT evaluates $\psi_m$
at the start of each sub-step (Euler explicit, Sec.~3), so the ratio of a load step lies between the
values of the rule at its start and its end; it does so to $3\times10^{-13}$. The figure plots it
against the mean stress ratio of each load step.

With the dilatancy cut-off ($e_0 = 0.60$, $e_{cv} = 0.65$), dilatancy stops when the void ratio
reaches $e_{cv}$. By Eq.~39 this happens at $\varepsilon_v = -\ln(1.65/1.60) = -3.08\,\%$; the
computed final value is also $-3.08\,\%$ (Fig.~\ref{fig:dilatancy}b; compare Fig.~4 of \citealp{Schanz_1999}).

\begin{figure}[h]
  \centering
  \begin{subfigure}[t]{0.48\textwidth}
    \centering
    \includegraphics[width=\textwidth]{dilatancy_rowe.pdf}
    \caption{mobilised dilatancy vs.\ friction angle}
  \end{subfigure}\hfill
  \begin{subfigure}[t]{0.48\textwidth}
    \centering
    \includegraphics[width=\textwidth]{dilatancy_ev.pdf}
    \caption{$\varepsilon_1$--$\varepsilon_v$ with and without cut-off}
  \end{subfigure}
  \caption{Drained triaxial compression, $\sigma_3 = 100$~kPa, $m = 0$, cap off.}
  \label{fig:dilatancy}
\end{figure}

% ---------------------------------------------------------------------------
\section{Cap: isotropic compression}

The cap is tested in normally consolidated isotropic compression with the unit-test parameters:
loading from $p = 100$ to $400$~kPa, unloading to $150$~kPa and reloading to $800$~kPa. Only the
cap is active on this path, and $\gamma^p$ stays zero. On primary loading
$dp/d\varepsilon_v = K_c\,((p + a)/(p\eref + a))^m$ with $K_c = K_s/(K_s/K_c)$ (Eqs.~31--32). On
unloading and reloading the modulus is $K_s$, with the same stress dependency.

Integrating these moduli gives the closed-form curve of Fig.~\ref{fig:cap}a. At the two reversal
points and at the end, the computed $\varepsilon_v$ deviates from it by 0.14\,\%, 0.44\,\% and
0.26\,\%; this is explicit integration error accumulated along the path. The tangent moduli deviate
by at most $1.5\times10^{-3}$ on primary loading and $2.8\times10^{-3}$ on unloading--reloading
(Fig.~\ref{fig:cap}b). Reloading becomes plastic again when $p$ reaches the pre-consolidation stress
$p_c = 402$~kPa reached in the first branch.

\begin{figure}[h]
  \centering
  \begin{subfigure}[t]{0.48\textwidth}
    \centering
    \includegraphics[width=\textwidth]{cap_isotropic.pdf}
    \caption{$\varepsilon_v$--$p$}
  \end{subfigure}\hfill
  \begin{subfigure}[t]{0.48\textwidth}
    \centering
    \includegraphics[width=\textwidth]{cap_bulk_modulus.pdf}
    \caption{tangent bulk modulus}
  \end{subfigure}
  \caption{Isotropic compression with an unloading--reloading loop (unit-test parameters).}
  \label{fig:cap}
\end{figure}

% ---------------------------------------------------------------------------
\section{General stress states}

Three drained paths start from an isotropic state of 100~kPa:
\begin{itemize}
  \item triaxial compression, $\sigma_2 = \sigma_3 = 100$~kPa;
  \item plane strain, $\varepsilon_{yy} = 0$ and $\sigma_{zz} = 100$~kPa, with the cap off as in the
        unit test;
  \item triaxial extension, $\sigma_{xx} = \sigma_{yy} = 100$~kPa with $\sigma_{zz}$ reduced.
\end{itemize}
All three reach the Mohr--Coulomb surface ($\sin\varphi_m = \sin\varphi$ to $10^{-8}$, failure flag
set). The final values of $b = (\sigma_2 - \sigma_3)/(\sigma_1 - \sigma_3)$ are 0, 0.25 and 1
(Fig.~\ref{fig:stress-states}). With $c = 0$, the Mohr--Coulomb surface normalised by $p$ is a fixed
hexagon in the deviatoric plane. The compression and extension paths end in its corners, and the
plane strain path ends on a face.

\begin{figure}[h]
  \centering
  \begin{subfigure}[t]{0.48\textwidth}
    \centering
    \includegraphics[width=\textwidth]{stress_states_mobilisation.pdf}
    \caption{mobilised friction vs.\ shear strain}
  \end{subfigure}\hfill
  \begin{subfigure}[t]{0.48\textwidth}
    \centering
    \includegraphics[width=\textwidth]{stress_states_pi_plane.pdf}
    \caption{deviatoric plane, stresses normalised by $p$}
  \end{subfigure}
  \caption{Drained compression, plane strain and extension from 100~kPa (unit-test parameters).}
  \label{fig:stress-states}
\end{figure}

% ---------------------------------------------------------------------------
\section{Hostun sand: calibration and oedometer test}
\label{sec:hostun-oedometer}

\citet{Schanz_1999} calibrate the model for loose Hostun sand (Sec.~6 and Table~1). Its cap input parameters,
$K_0^{NC}$ and $E_{oed}\eref$, are not inputs of this UMAT. Instead, $M$ and $K_s/K_c$ are found by
Newton iteration so that normally consolidated oedometer loading gives
\begin{itemize}
  \item a tangent ratio $K_0 = \Delta\sigma_3/\Delta\sigma_1 = 1 - \sin\varphi = 0.4408$
        (Sec.~5.2), and
  \item $E_{oed} = 16$~MPa at $\sigma_1 = p\eref$ (Eq.~37).
\end{itemize}
This gives $M = 1.233$ and $K_s/K_c = 2.109$. At the $K_0^{NC}$ stress ratio,
$\sin\varphi_m = 0.39 < \frac{3}{4}\sin\varphi$, so the shear surface has no plastic volume change
and the cap provides all of the plastic compaction.

Fig.~\ref{fig:oedometer} shows an oedometer test from $\sigma_1 = 2$~kPa (in a $K_0^{NC}$ state) to
200~kPa, unloading to 10~kPa and reloading to 300~kPa (compare Fig.~5 of \citealp{Schanz_1999}). Primary
loading follows Eq.~37 (Table~\ref{tab:oedometer}), with $K_0 = 0.441$ throughout, and $\varepsilon_1$
at 200~kPa is 1.821\,\% against 1.822\,\% from Eq.~37.

Unloading starts elastic, with $\Delta\sigma_3/\Delta\sigma_1 = \nu_{ur}/(1 - \nu_{ur}) = 0.25$. At
$\sigma_1 = 16$~kPa ($\sigma_3/\sigma_1 = 2.6$) the stress state reaches the shear hardening surface
in extension, and $\gamma^p$ increases again, from 0.0081 to 0.0106; this is the kink in
Fig.~\ref{fig:oedometer}b. The failure surface is not reached. Because the lateral strain is zero,
the plastic lateral expansion is balanced by elastic lateral compression, and $\sigma_3$ falls
faster than on the elastic line (28.5~kPa at $\sigma_1 = 10$~kPa instead of about 40~kPa).

Reloading is elastic with the same slope of 0.25, so it runs about 12~kPa below the unloading
line. The loop does not close, and the stress state stays more deviatoric than on primary loading
($\sigma_3/\sigma_1 = 0.404$ at $\sigma_1 = 169$~kPa). At this point it reaches the shear hardening
surface in compression, which was enlarged by the extension yielding because the hardening is
isotropic in $\gamma^p$. The cap is not reached first; $p_c$ only starts to increase again near
the previous maximum of $\sigma_1 = 200$~kPa. Plastic shear flow raises the tangent ratio
$\Delta\sigma_3/\Delta\sigma_1$ from 0.25 to about 0.42, which recovers part of the lost horizontal
stress. Above 200~kPa both surfaces are active and the path turns back towards the $K_0^{NC}$ line,
but only slowly, because the stiffness scales with stress. The tangent ratio stays just above the
secant ratio, so $\sigma_3/\sigma_1$ is 0.413 at 300~kPa, 0.425 at 1000~kPa and 0.432 at 3000~kPa.
This is why the reloading path in Fig.~\ref{fig:oedometer}b ends below the $K_0^{NC}$ line: the
stress state keeps a memory of the unloading--reloading loop.

\begin{table}[h]
  \centering
  \caption{Tangent oedometer modulus on primary loading, Hostun sand.}
  \label{tab:oedometer}
  \small
  \begin{tabular}{rccc}
    \toprule
    $\sigma_1$ [kPa] & $E_{oed}$ UMAT [MPa] & $E_{oed}$ Eq.~37 [MPa] & $K_0$ (tangent) \\
    \midrule
    50  & 10.20 & 10.19 & 0.441 \\
    100 & 16.01 & 16.00 & 0.441 \\
    200 & 25.12 & 25.11 & 0.441 \\
    \bottomrule
  \end{tabular}
\end{table}

\begin{figure}[h]
  \centering
  \begin{subfigure}[t]{0.48\textwidth}
    \centering
    \includegraphics[width=\textwidth]{hostun_oedometer.pdf}
    \caption{$\varepsilon_1$--$\sigma_1$}
  \end{subfigure}\hfill
  \begin{subfigure}[t]{0.48\textwidth}
    \centering
    \includegraphics[width=\textwidth]{hostun_oedometer_path.pdf}
    \caption{stress path $\sigma_1$--$\sigma_3$}
  \end{subfigure}
  \caption{Oedometer test on loose Hostun sand with an unloading--reloading loop.}
  \label{fig:oedometer}
\end{figure}

% ---------------------------------------------------------------------------
\section{Hostun sand: drained triaxial tests}
\label{sec:hostun-drained}

These are drained triaxial compression tests at $\sigma_3 = 300$~kPa (Fig.~6 of \citealp{Schanz_1999}) and
600~kPa to $\varepsilon_1 = 0.18$, followed by unloading to $q = 0$. At 300~kPa the stress ratio
reaches $(1 + \sin\varphi)/(1 - \sin\varphi) = 3.537$ at $\varepsilon_1 = 0.107$ and then stays
constant (Fig.~\ref{fig:drained}a); Fig.~6 of \citet{Schanz_1999} shows a maximum of about 3.5. With
$\psi = 0$ the shear surface has no plastic volume change. The soil compacts elastically and on the
cap until failure ($\varepsilon_v = 0.97\,\%$ at $\varepsilon_1 = 0.109$), and deforms at constant
volume afterwards (Fig.~\ref{fig:drained}b). Unloading is elastic.

The 600~kPa tests are the 300~kPa tests with all strains multiplied by $2^{1-m} = 1.275$, as in the
undrained tests (Sec.~\ref{sec:undrained}). With $\psi = 0$, failure is reached at
$\varepsilon_1 = 0.136$, with $\varepsilon_v = 1.24\,\%$.

The computed $\varepsilon_v$ in Fig.~6 of \citet{Schanz_1999} reaches about 0.9\,\% and then decreases after
failure. That decrease is dilatancy at failure, which $\psi = 0$ excludes (Eq.~13 then gives
$\varphi_{cv} = \varphi$, so $\psi_m = 0$ at failure).

With $\psi = 14^\circ$ (dashed lines), Eq.~13 gives $\varphi_{cv} = 21.5^\circ$, and the soil dilates
as soon as $\sin\varphi_m \ge \frac{3}{4}\sin\varphi$. At 300~kPa, $\varepsilon_v$ reaches its
maximum of 0.58\,\% at $\varepsilon_1 = 0.018$ and then decreases steadily, to $-8.3\,\%$ at the end
of loading. The stress ratio reaches the failure value earlier, at $\varepsilon_1 = 0.089$. At
600~kPa the maximum is 0.74\,\% at $\varepsilon_1 = 0.023$, failure is reached at
$\varepsilon_1 = 0.114$, and $\varepsilon_v = -7.5\,\%$ at the end of loading.

The secant stiffness at $q_f/2$ is 33.4~MPa, 18\,\% below
$E_{50} = E_{50}\eref\,(300/100)^{0.65} = 40.8$~MPa from Eq.~3 (Table~\ref{tab:e50}). This is caused
by the cap. The mean stress $p$ increases along the triaxial path from a normally consolidated
state, so the cap yields as well. Without the cap, Eq.~3 is recovered exactly. $\psi$ does not
affect the secant stiffness, since $\psi_m = 0$ below $\sin\varphi_m = \frac{3}{4}\sin\varphi$.

\begin{table}[h]
  \centering
  \caption{Secant stiffness at $q_f/2$ in drained triaxial compression of Hostun sand,
           $\sigma_3 = 300$~kPa (Eq.~3: 40.85~MPa).}
  \label{tab:e50}
  \small
  \begin{tabular}{lcc}
    \toprule
    Variant & $E_{50}$ secant [MPa] & $\varepsilon_v$ at $q_f/2$ [\%] \\
    \midrule
    no cap           & 40.85 & 0.19 \\
    cap (full model) & 33.45 & 0.49 \\
    \bottomrule
  \end{tabular}
\end{table}

\begin{figure}[h]
  \centering
  \begin{subfigure}[t]{0.48\textwidth}
    \centering
    \includegraphics[width=\textwidth]{hostun_drained_ratio.pdf}
    \caption{$\varepsilon_1$--$\sigma_1/\sigma_3$}
  \end{subfigure}\hfill
  \begin{subfigure}[t]{0.48\textwidth}
    \centering
    \includegraphics[width=\textwidth]{hostun_drained_ev.pdf}
    \caption{$\varepsilon_1$--$\varepsilon_v$}
  \end{subfigure}
  \caption{Drained triaxial tests on loose Hostun sand, $\sigma_3 = 300$ and 600~kPa: $\psi = 0$
           (Table~1 of \citealp{Schanz_1999}, solid) and $\psi = 14^\circ$ (dashed).}
  \label{fig:drained}
\end{figure}

% ---------------------------------------------------------------------------
\section{Hostun sand: undrained triaxial tests}
\label{sec:undrained}

Undrained (isochoric) triaxial compression runs from isotropic, normally consolidated states
$\sigma_c = 300$ and 600~kPa to $\varepsilon_1 = 0.20$ (compare Figs.~7--8 of \citealp{Schanz_1999}, whose
experiments are by \citet{Djedid_1986}). The total lateral stress stays at $\sigma_c$, so the excess
pore pressure is $\Delta u = \sigma_c - \sigma_3'$. Both $\psi = 0$ (Table~1 of \citealp{Schanz_1999}) and
$\psi = 14^\circ$ are computed (Table~\ref{tab:undrained}, Fig.~\ref{fig:undrained}).

Under undrained conditions $dp' = -K\,d\varepsilon_v^p$, so $p'$ changes only through plastic volume
change. With $\psi = 0$ the shear surface has none. $p'$ drops only through compaction on the cap,
which yields from the start because the samples are normally consolidated; without the cap, $p'$
stays at $\sigma_c$. The path rises steeply and bends to the left until it reaches the failure line
$q = 1.37\,p'$, at $\varepsilon_1 = 7.6\,\%$ for $\sigma_c = 300$~kPa and 9.7\,\% for 600~kPa. At
failure $\psi_m = 0$, so the state is stationary from there, at $\Delta u/\sigma_c = 0.555$ and
$\varphi'_{mob} = 34.0^\circ$.

With $\psi = 14^\circ$ the soil dilates from $\sin\varphi_m = \frac{3}{4}\sin\varphi$ on, well
before failure. $\Delta u/\sigma_c$ peaks at 0.412, at $\varepsilon_1 = 1.2\,\%$ for
$\sigma_c = 300$~kPa, and then decreases and becomes negative. $p'$ increases, and the path climbs
just below the failure line ($\varphi'_{mob} = 33.4^\circ$ at $\varepsilon_1 = 0.20$). Without a
dilatancy cut-off nothing limits this: at $\varepsilon_1 = 0.20$, $p' = 5.8$~MPa and
$\Delta u/\sigma_c = -9.6$ for $\sigma_c = 300$~kPa. In a real test, cavitation of the pore water
would limit the suction long before that. The stress ratio stays below that of the $\psi = 0$ test
after $\varepsilon_1 \approx 2\,\%$, because $q_f$ increases with $p'$ and the shear surface lags
behind it.

With $c = 0$ and stiffness laws that are power laws of stress, the model has no stress scale of its
own. The response normalised by $\sigma_c$ is therefore the same for every cell pressure, and only the
strains scale, with $\sigma_c^{1-m}$: the 600~kPa test is the 300~kPa test with all strains multiplied
by $2^{1-m} = 1.275$. The UMAT satisfies this to $6\times10^{-14}$ (checked for $\psi = 14^\circ$).

\subsection{Computed curves of \texorpdfstring{\citet{Schanz_1999}}{Schanz et al. (1999)}}
\label{sec:undrained-paper}

The computed undrained curves in Figs.~7--8 of \citet{Schanz_1999} cannot result from the parameters of
their Table~1:
\begin{itemize}
  \item \emph{Dilatancy with $\psi = 0$.} Beyond the failure line the computed paths of \citet{Schanz_1999}
        climb along it while $\Delta u$ decreases. With $\psi = 0$, Eq.~13 gives
        $\varphi_{cv} = \varphi$, so $\psi_m = 0$ at failure. There is no plastic volume change, $p'$
        is stationary and $\Delta u$ cannot decrease. The decrease requires dilatancy at failure,
        i.e.\ $\psi > 0$. The same holds for the computed drained test, whose $\varepsilon_v$
        decreases after failure (Sec.~\ref{sec:hostun-drained}).
  \item \emph{Dependence on the cell pressure.} By the scaling above, the 600~kPa test must reach a
        given $\Delta u/\sigma_c$ at a larger strain than the 300~kPa test, and at equal stress ratio
        $\sigma_1'/\sigma_3'$ both must have the same $\Delta u/\sigma_c$. In Fig.~8 of \citet{Schanz_1999} the
        600~kPa $\Delta u$ curve falls faster than the 300~kPa curve, and at equal stress ratio the
        two differ. This holds for any parameter set with $c = 0$, so no value of $\psi$ fits both
        tests. \citet{Schanz_1999} do not give enough information to find the cause; one possibility is that
        the line styles of the two $\Delta u$ curves are exchanged.
\end{itemize}

\begin{table}[h]
  \centering
  \caption{Undrained triaxial tests on loose Hostun sand (800 steps to $\varepsilon_1 = 0.20$).}
  \label{tab:undrained}
  \small
  \begin{tabular}{llcc}
    \toprule
    $\sigma_c$ & Quantity & $\psi = 0$ & $\psi = 14^\circ$ \\
    \midrule
    300~kPa & maximum $\Delta u/\sigma_c$                        & 0.555      & 0.412        \\
            & $\varepsilon_1$ at the maximum                     & 0.076      & 0.012        \\
            & $(p', q)$ at $\varepsilon_1 = 0.20$ [kPa]          & (246, 339) & (5768, 7787) \\
            & $\Delta u/\sigma_c$ at $\varepsilon_1 = 0.20$      & 0.555      & $-9.58$      \\
            & $\varphi'_{mob}$ at $\varepsilon_1 = 0.20$         & $34.0^\circ$ & $33.4^\circ$ \\
    \midrule
    600~kPa & maximum $\Delta u/\sigma_c$                        & 0.555      & 0.412        \\
            & $\varepsilon_1$ at the maximum                     & 0.097      & 0.016        \\
            & $(p', q)$ at $\varepsilon_1 = 0.20$ [kPa]          & (493, 677) & (6991, 9396) \\
            & $\Delta u/\sigma_c$ at $\varepsilon_1 = 0.20$      & 0.555      & $-5.43$      \\
            & $\varphi'_{mob}$ at $\varepsilon_1 = 0.20$         & $34.0^\circ$ & $33.3^\circ$ \\
    \bottomrule
  \end{tabular}
\end{table}

\begin{figure}[h]
  \centering
  \begin{subfigure}[t]{0.48\textwidth}
    \centering
    \includegraphics[width=\textwidth]{hostun_undrained_pq.pdf}
    \caption{$p'$--$q$}
  \end{subfigure}\hfill
  \begin{subfigure}[t]{0.48\textwidth}
    \centering
    \includegraphics[width=\textwidth]{hostun_undrained_ratio.pdf}
    \caption{$\varepsilon_1$--$\sigma_1'/\sigma_3'$}
  \end{subfigure}

  \vspace{1ex}
  \begin{subfigure}[t]{0.48\textwidth}
    \centering
    \includegraphics[width=\textwidth]{hostun_undrained_du.pdf}
    \caption{$\varepsilon_1$--$\Delta u/\sigma_c$}
  \end{subfigure}
  \caption{Undrained triaxial tests on loose Hostun sand, normally consolidated: $\psi = 0$
           (Table~1 of \citealp{Schanz_1999}, solid) and $\psi = 14^\circ$ (dashed); (a) on logarithmic axes. The
           solid lines of both cell pressures coincide in (c) once failure is reached. Compare
           Figs.~7--8 of \citet{Schanz_1999}.}
  \label{fig:undrained}
\end{figure}

% ---------------------------------------------------------------------------
\section{Numerical aspects}
\label{sec:numerics}

Within each load step the UMAT sub-steps so that the elastic stress increment per sub-step stays
below 1\,\% of $\sigma_3 + a$. The step-size study repeats the Hostun drained triaxial test with 10,
40 and 360 load steps and compares each with a 720-step reference (Table~\ref{tab:steps},
Fig.~\ref{fig:steps}). Even the coarsest run, at $\Delta\varepsilon_1 = 1.8\,\%$ per step, stays
within 1\,\% of $q_f$. For the undrained test ($\sigma_c = 300$~kPa), the end state ($p'$ and
$\Delta u/\sigma_c$) with 200, 800 and 3200 steps is the same to six digits for $\psi = 0$, and
within $4\times10^{-5}$ for $\psi = 14^\circ$.

Invariance under rotation of the loading frame, continuity of the stress update in the strain
increment, and loading at zero confinement and in tension are covered by the unit tests
(Sec.~2).

\begin{table}[h]
  \centering
  \caption{Step-size dependence, drained triaxial test on Hostun sand ($\sigma_3 = 300$~kPa,
           $\varepsilon_1 = 0.18$), against 720 steps.}
  \label{tab:steps}
  \small
  \begin{tabular}{rccc}
    \toprule
    Steps & $\Delta\varepsilon_1$ & $\max|\Delta q|/q_f$ & $\max|\Delta\varepsilon_v|$ \\
    \midrule
    10  & $1.8\times10^{-2}$ & $8.8\times10^{-3}$ & $1.4\times10^{-4}$ \\
    40  & $4.5\times10^{-3}$ & $1.6\times10^{-3}$ & $1.3\times10^{-5}$ \\
    360 & $5.0\times10^{-4}$ & $1.6\times10^{-5}$ & $1.3\times10^{-7}$ \\
    \bottomrule
  \end{tabular}
\end{table}

\begin{figure}[h]
  \centering
  \includegraphics[width=0.55\textwidth]{step_size.pdf}
  \caption{Drained triaxial test on Hostun sand with 10, 40 and 360 load steps.}
  \label{fig:steps}
\end{figure}

% ---------------------------------------------------------------------------
\section{Open points}
\label{sec:open}

\begin{enumerate}
  \item \emph{Computed curves of \citet{Schanz_1999} for loose Hostun sand.} Their Table~1 gives
        $\psi = 0$, but their computed drained and undrained curves dilate after failure, which
        $\psi = 0$ excludes; and their 600~kPa undrained $\Delta u$ curve contradicts the
        $\sigma_c$-scaling of the model (Sec.~\ref{sec:undrained-paper}). Results are therefore given
        for both $\psi = 0$ and $\psi = 14^\circ$.
  \item \emph{$E_{50}$ calibration.} With an active cap, the secant stiffness of a computed drained
        triaxial test is lower than the input $E_{50}$, by 18\,\% for Hostun sand. If
        $E_{50}\eref$ is taken directly from the secant stiffness of a triaxial test, the computed
        response of such soils is therefore softer than the test; $E_{50}\eref$ should then be
        fitted by simulating the test.
\end{enumerate}

\bibliographystyle{plainnat}
\bibliography{../refs}

\end{document}
